\section{MineCLIP: convertir el Internet Video en una Language-conditioned Reward}

La carencia más difícil de la sección anterior es que el creative goal no tiene una reward fácil de escribir a mano. La estrategia de MineCLIP es aprender una representación común de fragmentos de video y descripciones textuales: si el contenido del video corresponde mejor al texto objetivo, la similitud es mayor, y puede usarse como dense feedback para el aprendizaje de la policy. Un \term{reward model (modelo de recompensa)} es un modelo que asigna a una trajectory, una observation o un outcome una retroalimentación escalar; aproxima el objetivo humano, pero no es el objetivo mismo.

\subsection{Contrastive Learning e InfoNCE}

El \term{contrastive learning (aprendizaje contrastivo)} aprende representaciones acercando las muestras que corresponden entre sí y alejando las que no corresponden. MineCLIP codifica los fragmentos de video de Minecraft como $v_i$ y el texto correspondiente como $t_i$; durante el entrenamiento, el mismo par es positive y los demás textos o videos del lote funcionan como negative, de modo que los video--text pair con semántica coherente quedan más cerca en el embedding space.

\lecturefigure{slide-22-mineclip-contrastive-reward.jpg}{MineCLIP aprende la alineación video--lenguaje y produce una similitud semántica}{Diapositiva V022 recuperada de la grabación oficial; 00:16:51}

\begin{knowledgebox}{Lectura de la figura: distinguir primero la señal de entrenamiento de la reward en despliegue}
En la etapa de entrenamiento se usan pares de video clip y transcript, de modo que los pares que corresponden obtengan puntajes altos; en la etapa de despliegue, el segmento más reciente de observation del agent y el texto objetivo se envían al modelo para obtener la similarity reward. En la figura, un puntaje alto indica mayor cercanía semántica, no que la tarea esté necesariamente completada; el modelo puede preferir atajos visuales, el fondo o coocurrencias frecuentes, por lo que se requieren un outcome check y una auditoría humana.
\end{knowledgebox}

Una InfoNCE loss video-to-text simplificada es:
\[
\mathcal{L}_{\mathrm{NCE}}
=-\frac{1}{B}\sum_{i=1}^{B}
\log\frac{\exp(\operatorname{sim}(v_i,t_i)/\tau)}
{\sum_{j=1}^{B}\exp(\operatorname{sim}(v_i,t_j)/\tau)}.
\]
Donde $B$ es el tamaño del lote, $v_i$ es el embedding del $i$-ésimo video, $t_i$ es el embedding del texto correspondiente, $\operatorname{sim}$ suele ser la cosine similarity de vectores normalizados y $\tau$ es la temperature. La loss solo exige que los pares correspondientes estén relativamente más cerca; no garantiza que la similarity esté calibrada como una probabilidad real de éxito.

Para un texto objetivo $g$, la reward semántica que usa la policy puede escribirse así:
\[
r_t^{\text{MineCLIP}}=\operatorname{sim}\!\left(E_v(o_{t-k:t}),E_t(g)\right),
\qquad
J(\theta)=\mathbb{E}_{\pi_\theta}\left[\sum_t\gamma^t r_t^{\text{MineCLIP}}\right].
\]
Donde $E_v$ y $E_t$ son, respectivamente, los encoder de video y de texto, $o_{t-k:t}$ es la observation window de los $k$ pasos más recientes, $g$ es el objetivo en lenguaje, $r_t^{\text{MineCLIP}}$ es la learned reward y $J(\theta)$ es el policy objective. Cuanto más dense es la reward, más fácil resulta la exploración; pero la policy también amplifica activamente el reward bias.

\subsection{Entrenar una Language-conditioned Policy con una Learned Reward}

MineCLIP por sí solo no es un controlador completo. El sistema sigue necesitando una RL policy que elija la action a partir de la observation actual y que actualice sus parámetros con la learned reward. Por ello, el objetivo en lenguaje interviene en dos lugares: en uno define la dirección semántica de la reward; en el otro puede servir como policy condition, para que una misma red ejecute tareas distintas según la instruction.

\lecturefigure{slide-23-rl-with-mineclip.jpg}{La reward de MineCLIP impulsa el language-conditioned reinforcement learning}{Diapositiva V023 recuperada de la grabación oficial; 00:17:03}

\begin{lstlisting}[caption={Pseudocódigo conceptual del ciclo de entrenamiento de RL condicionado por MineCLIP}]
for goal in sampled_language_goals:
    obs = env.reset(goal)
    recent_frames = []
    while not env.done():
        action = policy.sample(obs, goal)
        next_obs = env.step(action)
        recent_frames.append(next_obs.frame)
        reward = mineclip.similarity(recent_frames[-k:], goal)
        replay.add(obs, action, reward, next_obs, goal)
        policy.update(replay)
        obs = next_obs
\end{lstlisting}

\begin{warningbox}{El reward hacking se agrava a medida que la policy se vuelve más capaz}
Si MineCLIP confunde cierto fondo, un movimiento de cámara o un objeto local con el éxito, el RL buscará ese shortcut de forma sistemática. Es indispensable registrar a la vez el programmatic state, la revisión humana de los videos, distintas seed, objetivos adversariales y la reward trajectory; «la reward sube» solo indica que la señal de entrenamiento mejoró, no equivale automáticamente a que se haya cumplido la intención humana.
\end{warningbox}

\teachervoice{El ponente describe MineCLIP como el puente entre el passive internet video y el active reinforcement learning: el video aporta la similitud semántica y luego el agent aprende a controlar mediante la interacción con el entorno. Este puente es importante, pero en la clase tampoco se afirma que la learned reward sea perfecta; sigue limitada por sesgos de los datos, ambigüedad de los objetivos y shortcut visuales.}

\subsection{Resultados: éxito en las tareas y robustez visual}

La diapositiva de resultados debe responder dos preguntas distintas: si la reward de MineCLIP basta para entrenar una policy que cumpla varios objetivos en lenguaje, y si la policy solo memoriza los escenarios de entrenamiento. Para la primera, la evidencia es el success de las tareas o la comparación con un baseline; para la segunda, la evidencia es la zero-shot visual condition. Ambas deben acotarse a las tareas evaluadas, las seed y la Minecraft distribution.

\lecturefigure{slide-24-mineclip-results.jpg}{Resultados del RL condicionado por MineCLIP en varios tipos de tareas}{Diapositiva V024 recuperada de la grabación oficial; 00:19:00}

\begin{knowledgebox}{Lectura de la figura: revisar primero la definición de las columnas y después la mejora relativa}
La tabla compara el indicador de success de distintas tareas y métodos. Primero hay que confirmar que la tarea de cada columna use el mismo evaluation protocol y luego comparar la reward de MineCLIP con las demás reward/baseline; la conclusión principal es que la representación video--lenguaje puede aportar una señal de entrenamiento útil, no que todas las tareas estén resueltas. El promedio oculta las tareas fallidas y la seed variance, y no permite inferir generality en un mundo abierto.
\end{knowledgebox}

\lecturefigure{slide-25-policy-robustness.jpg}{La policy aprendida conserva cierta robustez en condiciones visuales no vistas}{Diapositiva V025 recuperada de la grabación oficial; 00:19:57}

\begin{knowledgebox}{Lectura de la figura: la robustez proviene de la invariancia de la tarea, no de una capacidad OOD arbitraria}
La figura muestra distintos climas, iluminaciones, biomas o condiciones visuales. Primero conviene comparar si la conducta objetivo se mantiene y después verificar si esos cambios siguen perteneciendo a las reglas del mundo de entrenamiento. El resultado respalda que la policy no quedó totalmente ligada a una sola imagen, pero no demuestra que pueda transferirse a un juego nuevo, a un nuevo espacio de acciones ni a un robot real; la visual robustness y la embodiment transfer son niveles distintos.
\end{knowledgebox}

\teachervoice{La postura de la clase ante los resultados es que «la reward representation es útil», no que «Minecraft ya está resuelto». El ponente destaca la robustez bajo unseen visual condition porque se acerca más a la foundation-model transfer; pero también califica MineDojo, Voyager y Eureka como apenas un primer paso, scratching the surface.}

\subsection{Resumen de la sección}

MineCLIP usa el contrastive video--language learning para convertir objetivos en lenguaje natural en una dense reward, con la que después una RL policy aprende a controlar en el entorno. Demuestra que el internet video puede ser una fuente de supervisión para el active learning y, al mismo tiempo, expone el riesgo central de la learned reward: la similitud no equivale al éxito, la policy amplifica los sesgos del reward model y los resultados deben interpretarse junto con el estado programático y la auditoría humana.
